\section{Estimation, Bregman geometry, and identifiability}
\label{app:estimation}

\subsection{IC-LL is a Bregman divergence}

Write the interval counts $\sfont C=(\sfont C_i)$ and compensators
$\Xi(\theta)=(\Xi_i(\theta))$. A \emph{Bregman divergence} for a strictly convex
generator $\varphi$ is
$B_\varphi(x,y)=\varphi(x)-\varphi(y)-\langle x-y,\nabla\varphi(y)\rangle\ge0$.

\begin{proposition}\label{prop:kl}
With the generalized-KL generator $\varphi(x)=\sum_i(x_i\log x_i-x_i)$,
$\mathcal L_{\mathrm{IC}}(\theta)=\KL(\sfont C,\Xi(\theta))+\Gamma(\sfont C)$ with
$\Gamma$ free of $\theta$; hence
$\argmin_\theta\mathcal L_{\mathrm{IC}}=\argmin_\theta\KL(\sfont C,\Xi(\theta))$.
The squared-error generator $\tfrac12\|x\|^2$ instead yields the least-squares loss
$\sum_i(\sfont C_i-\Xi_i)^2$.
\end{proposition}

\begin{theorem}[Bregman $\leftrightarrow$ exponential family; Banerjee et al.]
For a regular exponential family with cumulant $\psi$ and mean $\mu=\nabla\psi$, the
negative log-likelihood equals $B_{\psi^\star}(x,\mu)$ up to a constant, where
$\psi^\star$ is the convex conjugate. Poisson gives $\psi^\star(\mu)=\mu\log\mu-\mu$
(so $B=\KL$); Gaussian gives $\psi^\star(\mu)=\mu^2/2$ (so $B=$ squared error).
\end{theorem}

\begin{intuition}
Maximum likelihood in any exponential family is ``move the model mean as close as
possible to the data, in the geometry set by the family's variance function.''
Choosing the divergence \emph{is} choosing the noise model: $x\log x\to$ Poisson
$\to$ KL $\to$ IC-LL; $\tfrac12x^2\to$ Gaussian $\to$ least squares. Because a
Poisson bucket has variance equal to its mean, KL correctly down-weights
large-count buckets by $1/\Xi_i$, whereas least squares treats all buckets alike
and is dominated by the loud ones --- so KL is the more robust loss when counts span
orders of magnitude (bursts, near criticality).
\end{intuition}

\subsection{Consistency, information, and goodness of fit}

\begin{theorem}[Asymptotics of the MLE; Ogata]
For a stationary, ergodic, subcritical process with true $\theta_0$ in the interior
of a compact, identifiable $\Theta$, twice-differentiable log-intensity with
integrable dominating envelopes, and finite nonsingular per-time information
$\mathfrak i(\theta_0)$, the MLE satisfies $\hat\theta_T\to\theta_0$ a.s.\ and
$\sqrt T(\hat\theta_T-\theta_0)\to\mathcal N(0,\mathfrak i(\theta_0)^{-1})$.
\end{theorem}
The information is
$\mathcal I(\theta)=\E\!\int_0^T\frac{\partial_\theta\lambda\,\partial_\theta\lambda^\top}{\lambda}\dd t$
(event-time) or
$\mathcal I^{\mathrm{IC}}_{jk}=\sum_i\partial_{\theta_j}\Xi_i\,\partial_{\theta_k}\Xi_i/\Xi_i$
(interval-censored); standard errors come from inverting the observed information
(the Hessian of $-\ell$), and Cram\'er--Rao says no unbiased estimator beats
$\mathcal I^{-1}$.

\begin{theorem}[Time-rescaling; goodness of fit]
If $\Lambda$ is continuous, strictly increasing with $\Lambda(\infty)=\infty$, the
rescaled times $\Lambda(t_k)$ form a unit-rate Poisson process; equivalently the
rescaled gaps $\Lambda(t_k)-\Lambda(t_{k-1})$ are i.i.d.\ $\mathrm{Exp}(1)$.
\end{theorem}

\begin{intuition}
Measure time not in seconds but in \emph{accumulated compensator} --- ``operational
time'' that ticks fast when the intensity is high. In that clock the events look
like a plain unit-rate Poisson process, giving a model-free residual check
(rescaled gaps should be $\mathrm{Exp}(1)$).
\end{intuition}

\subsection{The $\kappa$--$\theta$ ridge, and the cost of covariates}

\begin{proposition}[Weak identifiability of the timescale]
In $\phi=\alpha e^{-\beta t}$, $\alpha=\kappa\theta$, $\beta=\theta$, the
$(\alpha,\beta)$ information matrix is typically ill-conditioned, with its
small-eigenvalue direction aligned with curves of constant $\alpha/\beta=\kappa$.
Hence $\kappa$ is well identified while $\theta$ is weak, and the more so the
coarser the censoring: as the bucket width grows the timing-borne sensitivity
$\partial_\theta\Xi_i\to0$ while $\partial_\kappa\Xi_i$ stays order one.
\end{proposition}

\begin{intuition}
The counts pin down ``how much'' (the gain $\kappa$) but not ``how fast'' (the
timescale $\theta$), because a bucket integrates away the within-bucket timing that
carries $\theta$. So treat $\kappa$ and the covariate effects $\gamma$ as the
trustworthy estimands; fix or prior-constrain $\theta$ from domain knowledge and
report it with wide intervals.
\end{intuition}

\paragraph{Excitation covariates add non-convexity.} Baseline covariates
(\S\ref{sec:baseline-cov}) enter the compensator linearly through a log-link in the
\emph{forcing} and keep the problem well-behaved; the covariate effect $\gamma$ is
recovered tightly. Excitation covariates (\S\ref{sec:excitation-cov}) enter through
$e^{\delta^\top Z}$ \emph{inside} the intensity dynamics, so the likelihood is
non-convex in $\delta$, with up to $pM^2$ extra parameters and a real need for
covariate--event co-occurrence to identify each $\delta_{mj}$. This is the
statistical half of ``excitation covariates are harder''; the computational half is
that $\Xi$ now comes from integrating an LTV ODE rather than from a formula.
